\section*{Abstract}
Deep reinforcement learning has achieved remarkable success in domains such as game playing, robotics, recommendation systems, control, and large-scale decision making \citep{mnih2015human,silver2017mastering}. However, despite these successes, modern deep reinforcement learning systems frequently suffer from a phenomenon known as plasticity loss.

Plasticity in general is the ability of a system to change how it perceives and processes the data and be able to adapt. Considering the human brain and the studies performed on it, neuroplasticity is the brain's remarkable capacity to reorganize, grow, and rewire its neural connections throughout an individual’s lifetime. Rather than being a fixed organ, it continuously adapts to new experiences, learning processes, and environmental stimuli, or following an injury \citep{puderbaugh2023neuroplasticity}. In the context of deep learning and neural networks, it refers to a neural network’s ability to adapt its predictions when new information becomes available. As training progresses, networks often become increasingly difficult to modify, resulting in slower learning, reduced adaptability, and performance degradation in non-stationary environments \citep{lyle2023understanding,ash2020warm}.

In recent years, we have seen an increasing rate in the adoption of RL systems in various domains, from self-driving cars to autonomous robots and superhuman results on many games. As we use DRL systems more and more by the day, we need to find the root causes preventing neural networks from benefitting of plasticity while solving tasks in the reinforcement learning regime.

The objective of this project is to conduct a literature review and investigate the mechanisms responsible for plasticity loss and understand how this phenomenon affects deep reinforcement learning systems. We base our benchmark on the reproduction of the results from \citep{lyle2023understanding} and argue about the effect of geometrical landscape of functions and their effect on learning.
